
        \documentclass[fleqn]{article}      
        \usepackage{amsmath}
        \usepackage{fontspec}
        \usepackage{kotex} % 한국어 지원  

        \begin{document}      
        ```latex
\documentclass{article}
\usepackage{amsmath}
\usepackage{geometry}
\geometry{a4paper, margin=1in}
\begin{document}

\section*{고등수학 (상) - 방정식과 부등식 - 연립이차방정식}

\subsection*{문제 1 (쉬움, 연립이차방정식)}
다음 연립이차방정식을 푸시오.
\[
\begin{cases}
x^2 + y^2 = 25 \\
x + y = 7
\end{cases}
\]
\begin{enumerate}
    \item (3,4) 
    \item (4,3) 
    \item (5,2) 
    \item (6,1) 
    \item (7,0)
\end{enumerate}

\subsection*{문제 2 (쉬움, 연립이차방정식)}
다음 연립이차방정식을 푸시오.
\[
\begin{cases}
x^2 + y^2 = 10 \\
x - y = 2
\end{cases}
\]
\begin{enumerate}
    \item (3,1) 
    \item (1,3) 
    \item (2,2) 
    \item (4,0) 
    \item (0,4)
\end{enumerate}

\subsection*{문제 3 (쉬움, 연립이차방정식)}
다음 연립이차방정식을 푸시오.
\[
\begin{cases}
x^2 + y^2 = 16 \\
x + y = 4
\end{cases}
\]
\begin{enumerate}
    \item (0,4) 
    \item (2,2) 
    \item (4,0) 
    \item (3,1) 
    \item (1,3)
\end{enumerate}

\subsection*{문제 4 (쉬움, 연립이차방정식)}
다음 연립이차방정식을 푸시오.
\[
\begin{cases}
x^2 + y^2 = 20 \\
x - y = 4
\end{cases}
\]
\begin{enumerate}
    \item (6,2) 
    \item (2,6) 
    \item (5,3) 
    \item (3,5) 
    \item (4,4)
\end{enumerate}

\subsection*{문제 5 (쉬움, 연립이차방정식)}
다음 연립이차방정식을 푸시오.
\[
\begin{cases}
x^2 + y^2 = 18 \\
x + y = 6
\end{cases}
\]
\begin{enumerate}
    \item (3,3) 
    \item (4,2) 
    \item (2,4) 
    \item (6,0) 
    \item (0,6)
\end{enumerate}

\subsection*{문제 6 (중간, 연립이차방정식)}
연립이차방정식
\[
\begin{cases}
x^2 + y^2 = 13 \\
xy = 6
\end{cases}
\]
의 해를 구하시오.
\begin{enumerate}
    \item (3,2) 
    \item (2,3) 
    \item (4,1) 
    \item (1,4) 
    \item (5,0)
\end{enumerate}

\subsection*{문제 7 (중간, 연립이차방정식)}
연립이차방정식
\[
\begin{cases}
x^2 + y^2 = 50 \\
x - y = 5
\end{cases}
\]
의 해를 구하시오.
\begin{enumerate}
    \item (5,5) 
    \item (10,0) 
    \item (7,3) 
    \item (8,2) 
    \item (6,4)
\end{enumerate}

\subsection*{문제 8 (중간, 연립이차방정식)}
연립이차방정식
\[
\begin{cases}
x^2 + y^2 = 29 \\
x + y = 8
\end{cases}
\]
의 해를 구하시오.
\begin{enumerate}
    \item (5,3) 
    \item (3,5) 
    \item (4,4) 
    \item (6,2) 
    \item (2,6)
\end{enumerate}

\subsection*{문제 9 (중간, 연립이차방정식)}
연립이차방정식
\[
\begin{cases}
x^2 + y^2 = 45 \\
xy = 14
\end{cases}
\]
의 해를 구하시오.
\begin{enumerate}
    \item (5,5) 
    \item (3,6) 
    \item (7,2) 
    \item (2,7) 
    \item (4,5)
\end{enumerate}

\subsection*{문제 10 (중간, 연립이차방정식)}
연립이차방정식
\[
\begin{cases}
x^2 + y^2 = 34 \\
x - y = 4
\end{cases}
\]
의 해를 구하시오.
\begin{enumerate}
    \item (8,4) 
    \item (6,2) 
    \item (4,0) 
    \item (5,3) 
    \item (3,5)
\end{enumerate}

\subsection*{문제 11 (어려움, 연립이차방정식)}
연립이차방정식
\[
\begin{cases}
x^2 - y^2 = 12 \\
xy = 16
\end{cases}
\]
을 만족하는 \( (x, y) \)의 한 쌍을 구하시오.

\subsection*{문제 12 (어려움, 연립이차방정식)}
연립이차방정식
\[
\begin{cases}
x^2 + y^2 = 41 \\
x + y = 9
\end{cases}
\]
을 만족하는 \( (x, y) \)의 한 쌍을 구하시오.

\subsection*{문제 13 (어려움, 연립이차방정식)}
연립이차방정식
\[
\begin{cases}
x^2 + y^2 = 30 \\
xy = 8
\end{cases}
\]
을 만족하는 \( (x, y) \)의 한 쌍을 구하시오.

\subsection*{문제 14 (어려움, 연립이차방정식)}
연립이차방정식
\[
\begin{cases}
x^2 + y^2 = 20 \\
x - y = 3
\end{cases}
\]
을 만족하는 \( (x, y) \)의 한 쌍을 구하시오.

\subsection*{문제 15 (어려움, 연립이차방정식)}
연립이차방정식
\[
\begin{cases}
x^2 + y^2 = 25 \\
xy = 12
\end{cases}
\]
을 만족하는 \( (x, y) \)의 한 쌍을 구하시오.

\section*{정답 및 해설}

\subsection*{문제 1 해설}
주어진 연립이차방정식을 대입법을 사용하여 풀면,  
\[
(x+y)^2 - 2xy = 25
\]
\[
49 - 2xy = 25
\]
\[
xy = 12
\]
따라서 정답은 \( \boxed{(3,4)} \) 또는 \( \boxed{(4,3)} \)이다.

\subsection*{문제 2 해설}
주어진 연립이차방정식을 대입법을 사용하여 풀면,  
\[
x = y + 2 \implies (y + 2)^2 + y^2 = 10
\]
\[
y^2 + 4y + 4 + y^2 = 10 \implies 2y^2 + 4y - 6 = 0 \implies y = 1 \text{ 또는 } y = -3
\]
따라서 정답은 \( \boxed{(3,1)} \) 또는 \( \boxed{(-1,-3)} \)이다.

\subsection*{문제 3 해설}
주어진 연립이차방정식을 대입법을 사용하여 풀면,  
\[
(x+y)^2 - 2xy = 16
\]
\[
16 - 2xy = 16 \implies xy = 0
\]
따라서 정답은 \( \boxed{(4,0)} \) 또는 \( \boxed{(0,4)} \)이다.

\subsection*{문제 4 해설}
주어진 연립이차방정식을 대입법을 사용하여 풀면,  
\[
x = y + 4 \implies (y + 4)^2 + y^2 = 20
\]
\[
y^2 + 8y + 16 + y^2 = 20 \implies 2y^2 + 8y - 4 = 0 \implies y = 1 \text{ 또는 } y = -2
\]
따라서 정답은 \( \boxed{(6,2)} \) 또는 \( \boxed{(2,6)} \)이다.

\subsection*{문제 5 해설}
주어진 연립이차방정식을 대입법을 사용하여 풀면,  
\[
x + y = 6 \implies y = 6 - x
\]
\[
x^2 + (6-x)^2 = 18 \implies x^2 + 36 - 12x + x^2 = 18 \implies 2x^2 - 12x + 18 = 0 \implies x = 3 \text{ 또는 } x = 1.5
\]
따라서 정답은 \( \boxed{(3,3)} \) 또는 \( \boxed{(1.5,4.5)} \)이다.

\subsection*{문제 6 해설}
이 문제는 대입법으로 풀 수 있다.  
\[
xy = 6 \implies y = \frac{6}{x}
\]
\[
x^2 + \left(\frac{6}{x}\right)^2 = 13 \implies x^4 - 13x^2 + 36 = 0
\]
해를 구하면 \( x = 3, y = 2 \) 또는 \( x = 2, y = 3 \)이다.

\subsection*{문제 7 해설}
이 문제는 대입법으로 풀 수 있다.  
\[
x - y = 5 \implies x = y + 5
\]
\[
(y + 5)^2 + y^2 = 50 \implies 2y^2 + 10y - 25 = 0
\]
해를 구하면 \( x = 10, y = 5 \) 또는 \( x = 5, y = 0 \)이다.

\subsection*{문제 8 해설}
대입법으로 풀면,  
\[
x + y = 8 \implies y = 8 - x
\]
\[
x^2 + (8 - x)^2 = 29 \implies 2x^2 - 16x + 35 = 0
\]
해를 구하면 \( x = 5, y = 3 \) 또는 \( x = 3, y = 5 \)이다.

\subsection*{문제 9 해설}
대입법으로 풀면,  
\[
xy = 14 \implies y = \frac{14}{x}
\]
\[
x^2 + \left(\frac{14}{x}\right)^2 = 45 \implies x^4 - 45x^2 + 196 = 0
\]
해를 구하면 \( x = 7, y = 2 \) 또는 \( x = 2, y = 7 \)이다.

\subsection*{문제 10 해설}
대입법으로 풀면,  
\[
x - y = 4 \implies x = y + 4
\]
\[
(y + 4)^2 + y^2 = 34 \implies 2y^2 + 8y - 18 = 0
\]
해를 구하면 \( x = 8, y = 4 \) 또는 \( x = 4, y = 0 \)이다.

\subsection*{문제 11 해설}
이 문제는 대입법으로 풀 수 있다.  
\[
x^2 - y^2 = 12 \implies (x-y)(x+y) = 12
\]
\[
xy = 16 \implies y = \frac{16}{x}
\]
결과적으로 \( (x,y) \)의 해는 \( (8,2) \), \( (4,-4) \) 등이 있다.

\subsection*{문제 12 해설}
대입법으로 풀면,  
\[
x + y = 9 \implies y = 9 - x
\]
\[
x^2 + (9 - x)^2 = 41 \implies 2x^2 - 18x + 18 = 0
\]
해를 구하면 \( x = 6, y = 3 \) 또는 \( x = 3, y = 6 \)이다.

\subsection*{문제 13 해설}
대입법으로 풀면,  
\[
xy = 8 \implies y = \frac{8}{x}
\]
\[
x^2 + \left(\frac{8}{x}\right)^2 = 30 \implies x^4 - 30x^2 + 64 = 0
\]
해를 구하면 \( x = 4, y = 2 \) 또는 \( x = 2, y = 4 \)이다.

\subsection*{문제 14 해설}
대입법으로 풀면,  
\[
x - y = 3 \implies x = y + 3
\]
\[
(y + 3)^2 + y^2 = 20 \implies 2y^2 + 6y - 11 = 0
\]
해를 구하면 \( x = 7, y = 4 \) 또는 \( x = 4, y = 1 \)이다.

\subsection*{문제 15 해설}
대입법으로 풀면,  
\[
xy = 12 \implies y = \frac{12}{x}
\]
\[
x^2 + \left(\frac{12}{x}\right)^2 = 25 \implies x^4 - 25x^2 + 144 = 0
\]
해를 구하면 \( x = 6, y = 2 \) 또는 \( x = 2, y = 6 \)이다.

\end{document}
```      
        \end{document} 
    